\chapter{Ontwerp}%
\label{ch:ontwerp}

Het ontwerp van de XR-applicatie vormt de fundamentele brug tussen de medische handelingen en de technologische realisatie op de Meta Quest 3. Om de onderzoeksvraag te beantwoorden met betrekking tot de technische ondersteuning van medische simulaties, is gekozen voor een schaalbaar en modulair systeemontwerp dat de focus legt op een flexibele software-architectuur en didactische bruikbaarheid. Het ontwerp overstijgt een loutere visuele representatie door een uitbreidbare logische structuur te creëren waarop in de toekomst eenvoudig nieuwe klinische scenario's kunnen worden aangesloten binnen de beoogde 'HOGENT XR Academy' \cite{Radianti2020VRHE}.

\section{Modulaire XR-trainingsarchitectuur en scenariomanagement}

De interne scènestructuur is zodanig ontworpen dat de visuele assets en de logische interactiecomponenten strikt gescheiden blijven. Dit bevordert de schaalbaarheid van het platform, waardoor het mogelijk wordt om in de toekomst met minimale aanpassingen nieuwe klinische omgevingen en handelingen te laden zonder de kernlogica van het basissysteem te verstoren. Deze modulaire opbouw is essentieel voor de gewenste uitbreidbaarheid naar een bredere HOGENT XR Academy \cite{Radianti2020VRHE}.

Binnen de gerealiseerde PoC is de effectieve runtime-scope echter afgebakend tot twee operationele scenes: TrainingLobby en MedicalTrainingScene. Andere scenes in het project zijn in deze iteratie niet meegenomen in de evaluatie van functionaliteit of leerlogica.

De kern van de applicatie wordt aangestuurd door een centrale TrainingManager. Deze manager fungeert als een Singleton en is verantwoordelijk voor het initialiseren en inladen van een specifieke trainingsmodule. Om de uitbreidbaarheid te garanderen, erven alle medische scenario's over van een generieke basisklasse genaamd TrainingModuleBase. 

Binnen een actieve module wordt de procedurele flow beheerd door een ScenarioController. Deze controller doorloopt opeenvolgende scenario-states om de stappen van de gebruiker logisch te sturen. Voor deze PoC is één specifieke module uitgewerkt, namelijk de AED\_CPR\_Module. De overgangen tussen de statussen binnen deze module worden getriggerd door een event-driven systeem. Wanneer een gebruiker een interactie voltooit, zoals het voltooien van een reanimatiecyclus, vuurt de actieve interactor een event af, waarna de ScenarioController de status valideert en de transitie naar de volgende fase inzet.

\begin{figure}[htbp]
    \centering
    \includegraphics[width=0.9\textwidth]{architecture.jpg}
    \caption{Architectuur van de PoC (eigen ontwerp).}
    \label{fig:architure}
\end{figure}

\section{Technisch ontwerp van de interactoren}

Om een realistische trainingsomgeving te creëren, is de interactielaag ontworpen als een abstractielaag tussen de hardware-input van de Meta Quest 3 en de scenario-logica. Er is gekozen voor een architectuur die gebruikmaakt van specifieke interactors voor verschillende medische handelingen, waarbij de nadruk ligt op fysieke accuraatheid en haptische validatie.

\subsection{Ruimtelijke logica en elektrodenvalidatie}

Voor de interactie met de AED-elektroden is een bruikbaar plaatsingssysteem geïmplementeerd binnen de virtuele omgeving om te garanderen dat de student de elektroden op de voorgeschreven posities positioneert. De technische registratie van de elektrodenplaatsing verloopt via specifieke trigger-zones op de patiënt-mesh die corresponderen met de anatomische locaties voor de apex en het sternum. Elke elektrode is voorzien van een standaard Grab Interactor uit de Unity XR Interaction Toolkit waarmee het object kan worden opgepakt en bewogen.

Wanneer een elektrode de vooraf gedefinieerde radius van de desbetreffende trigger-zone op de torso betreedt, registreert het systeem de fysieke aanwezigheid van het object. Bij een succesvolle plaatsing vuurt de interactielaag een event af naar de ScenarioController om te bevestigen dat de desbetreffende zone is geactiveerd. Zodra beide plaatsingsevents door de controller zijn opgevangen en gevalideerd, wordt deze stap binnen de scenarioflow als succesvol beschouwd en wordt de transitie naar de analysefase ingezet.

Om het inherente gebrek aan fysieke tegendruk en tactiele sensaties bij standalone XR-simulaties op te vangen, maakt de applicatie gebruik van haptische feedback als een vorm van pseudo-haptiek. Bij het uitvoeren van reanimatiecompressies en bij specifieke interacties in de handelingflow stuurt de software op het moment van detectie een trilling naar de controllers van de Meta Quest 3. Deze directe impuls biedt de gebruiker een onmiddellijke fysieke bevestiging dat de handeling door de software-architectuur is geregistreerd, wat helpt om de cognitieve belasting tijdens de training te verlagen.

\section{Functionele implementatie van de AED/CPR-scenarioflow}

De scenarioflow in de actieve PoC is geconcentreerd binnen de AED\_CPR\_Module. Deze module vormt de directe technische vertaling van het klinische reanimatie-algoritme naar de opeenvolgende toestanden binnen de ScenarioController \cite{ERC2025BLS}. In plaats van een allesomvattende simulatie van de volledige Basic Life Support-procedure, richt deze implementatie zich specifiek op de kritieke technische en procedurele interacties rondom hartmassage en defibrillatie. De module loodst de gebruiker opeenvolgend door vier functionele fasen om een gestructureerd verloop te garanderen.

Het scenario start met de assessmentfase waarin de student de eerste status van de patiënt beoordeelt en de interactie opstart. Hierna volgt de reanimatiefase waarin de feitelijke borstcompressies moeten worden uitgevoerd. Tijdens deze status analyseert de CPRInteractor continu de neerwaartse controllerbewegingen om te verifiëren of het tempo aansluit bij de gestelde marges. Wanneer aan de numerieke succesvoorwaarden van de reanimatiecyclus is voldaan, initieert het systeem automatisch de transitie naar de volgende fase.

De daaropvolgende fase betreft de AED-setup waarin het virtuele AED-toestel actief wordt ingezet binnen de simulatie-logica. De student dient de elektroden uit de koffer te nemen en deze op de daarvoor voorziene trigger-zones op de patiënt-mesh te plaatsen. Zodra de plaatsingsevents succesvol zijn opgevangen door de ScenarioController, start de afrondende schokfase. Het virtuele apparaat voert hierbij een gesimuleerde hartritmeanalyse uit en geeft vervolgens via gesproken instructies en ruimtelijke audio de schokknop vrij, mits dit medisch geïndiceerd is \cite{ERC2025BLS, ILCOR2023Immersive}.

\begin{figure}[htbp]
    \centering
    \includegraphics[width=0.9\textwidth]{ScenarioFlowState.jpg}
    \caption{State transitie diagram in de AED\_CPR\_Module aangedreven door de ScenarioController (eigen ontwerp).}
    \label{fig:flow_state}
\end{figure}

\section{UX/UI}

Voor de interface is gekozen voor een diegetisch ontwerp, wat betekent dat alle menu-elementen en feedbackonderdelen visueel zijn geïntegreerd in de virtuele wereld zelf. Er is bewust voor gekozen om geen platte gebruikersinterfaces aan de camera van de Meta Quest 3 te koppelen, aangezien dergelijke zwevende schermen de immersie verstoren en een bekende oorzaak zijn van simulator sickness \cite{Kennedy1993SSQ}. Informatie en statuswijzigingen verschijnen daardoor direct op de medische objecten, bijvoorbeeld via een subtiele visuele indicator op de behuizing van het AED-toestel.

Wanneer een student een procedurele fout maakt tijdens het scenario, zoals het aanhouden van een incorrect compressieritme, wordt de simulatie niet direct afgebroken. De applicatie registreert de afwijking en biedt de student de kans om de handeling binnen de actieve state te corrigeren, wat aansluit bij de onderwijskundige principes van deliberate practice \cite{Fowler2015VR}. Dit zorgt ervoor dat het systeem bruikbaar blijft als formatief leermiddel zonder de softwarematige flow te blokkeren.

De evaluatie en datacaptatie worden afgehandeld door een centrale EvaluationManager. Tijdens de simulatie registreert deze manager de unieke module-identifier, de totale verstreken sessietijd en de opgelopen fouten in de uitvoering. Deze samenvattende sessiegegevens worden lokaal persistent opgeslagen via PlayerPrefs. Parallel daarop registreert de LoggingManager gebeurtenisdata in een xAPI-geinspireerd Actor-Verb-Object-formaat en exporteert deze gegevens asynchroon naar een JSON-bestand in de persistente opslag van het toestel. Deze combinatie van lokale sessiesamenvatting en eventgebaseerde logging onderbouwt de technische haalbaarheid van objectieve en geautomatiseerde resultatenregistratie binnen het ontwikkelde platform.